\newpage

\addcontentsline{toc}{chapter}{\textbf{ABSTRACT}}
\setcounter{tocdepth}{0}

\begin{center}
    \Large \textbf{ABSTRACT}
\end{center}

Concept drift is a major cause of performance degradation in machine learning systems when operating data changes over time. This thesis proposes \textbf{SE-CDT (ShapeDD-Enhanced Concept Drift Type)}, a unified detector-classifier system for unsupervised drift handling on streaming data, paired with a type-specific adaptation framework.

The \textbf{detection module} of SE-CDT, named \textbf{ShapeDD-IDW}, extends ShapeDD with a two-stage design. A Standard MMD signal first locates drift candidates; candidate peaks are then validated by IDW-MMD (Inverse Density-Weighted MMD) using a Gamma-approximated null distribution, with Bonferroni correction when several peaks are tested in the same trace. On a 14-dataset benchmark with 30 independent runs, $\delta = 75$ tolerance, and $\text{cooldown} = 150$, the detection module reaches $F_1 = 0.531$, statistically tied with DAWIDD under the Nemenyi rank test and above the original ShapeDD in the same benchmark, while running approximately $7\times$ faster than ShapeDD because the Gamma null replaces the costly permutation test.

The \textbf{classification module} of SE-CDT identifies sudden, gradual, incremental, recurrent, and blip drift from the shape of the MMD signal without labels. It uses concept memory for recurrent drift and a relaxed blip rule that accepts two- or three-peak transient signals. The classifier reaches 80.1\% TCD/PCD category accuracy, 50.5\% micro-averaged subtype accuracy, and 50.0\% macro-averaged subtype accuracy. Sudden, blip, and recurrent drift are identified more reliably than gradual and incremental drift, with incremental drift remaining the weakest case.

For adaptation, type-specific model updates driven by SE-CDT's classification output improve prequential accuracy on Stepping and Sudden scenarios compared with no adaptation, reaching 74.27\% and 73.71\%, respectively. On the Mixed scenario, the method reaches 85.73\%, close to Periodic Retrain but with fewer redundant alarms. A Kafka/Redpanda-compatible prototype is also tested to validate the end-to-end flow from SE-CDT detection to model update and reload.
